\documentclass[conference]{IEEEtran}
\usepackage[T1]{fontenc}
\usepackage{amsmath}
\usepackage{booktabs}
\usepackage{cite}
\usepackage{capt-of}
\usepackage{tikz}
\usepackage[hidelinks]{hyperref}
\urlstyle{same}
\hypersetup{
  pdftitle={Measuring the CPU Frontier for Streaming Speech Recognition},
  pdfsubject={Word Error Rate, Inference Cost, and Sustained Stream Density},
  pdfauthor={Anonymous Authors},
}
\input{metrics.tex}

\title{Measuring the CPU Frontier for Streaming Speech Recognition\\
\large Word Error Rate, Inference Cost, and Sustained Stream Density}
\author{\IEEEauthorblockN{Anonymous Authors}
\IEEEauthorblockA{Preliminary results from the current checkpoints}}

\newcommand{\cps}{\mbox{CPU-s/audio-s}}
\newcommand{\metric}[2]{\paperhasvalue{#1}{#2}{\papernum{#1}{#2}}{\textemdash}}
\newcommand{\pctmetric}[2]{\paperhasvalue{#1}{#2}{\papernum{#1}{#2}\%}{\textemdash}}
\newcommand{\trainrow}[2]{#1 & \metric{#2}{wer_clean_pct} & \metric{#2}{wer_other_pct} \\}
\newcommand{\costrow}[2]{#1 & \metric{#2}{cpu_seconds_per_audio_second} & \metric{#2}{wer_pooled_pct} &
\metric{#2}{wer_clean_pct} & \metric{#2}{wer_other_pct} &
\metric{#2}{peak_rss_mib} & \metric{#2}{artifact_mib} \\}
\newcommand{\densityrow}[2]{#1 & \metric{#2}{streams_ccx} & \metric{#2}{streams_host} &
\metric{#2}{test_clean_wer_pct} & \metric{#2}{test_other_wer_pct} &
\metric{#2}{test_clean_delta_pp} & \metric{#2}{test_other_delta_pp} \\}
\begin{document}
\maketitle

\begin{abstract}
CPU speech services must process concurrent streams as audio arrives and
bound the age of undecoded audio. We evaluate a streaming transducer with
multirate encoding and prepared INT8 inference. The benchmark pairs pooled
word error rate (WER) with launch-through-exit CPU cost on a 730-utterance
LibriSpeech development panel and measures paced stream density separately.
In a rehearsal on one four-core Zen 2 complex of a Threadripper Pro 3995WX,
four workers cycled ten dev-clean recordings per stream in 20-second rungs
under an 800-ms per-worker p99 lateness deadline. Prepared INT8 D2MR at step 50,000
passed at \papernum{rehearsal-mr-pass}{streams_ccx} streams versus
\papernum{rehearsal-r0-pass}{streams_ccx} for Zipformer INT8
(\papernum{rehearsal-mr-pass}{streams_per_core} versus
\papernum{rehearsal-r0-pass}{streams_per_core} per core), a
\papernum{rehearsal-mr-pass}{pass_rung_ratio_to_r0}-fold difference in
passing rungs. Worst-worker p99 lateness was
\papernum{rehearsal-mr-pass}{p99_lateness_ms} and
\papernum{rehearsal-r0-pass}{p99_lateness_ms} ms, respectively.
At the common \papernum{rehearsal-r0-common}{streams_ccx}-stream rung, D2MR used
\papernum{rehearsal-mr-common}{cpu_seconds_per_audio_second} CPU seconds per
audio second versus \papernum{rehearsal-r0-common}{cpu_seconds_per_audio_second}
for Zipformer. On the 730-utterance development panel, the same configurations
achieved \pctmetric{rehearsal-mr-score}{wer_pooled_pct} and
\pctmetric{rehearsal-r0-score}{wer_pooled_pct} pooled WER, a gap of
\papernum{rehearsal-mr-score}{wer_pooled_delta_pp} percentage points
(descriptive 95\% interval,
\papernum{rehearsal-mr-score}{wer_pooled_lower_pp}--\papernum{rehearsal-mr-score}{wer_pooled_upper_pp}).
Matched cold cost, full-host capacity, and registered test WER remain pending.
\end{abstract}

\begin{IEEEkeywords}
streaming speech recognition, CPU inference, word error rate, quantization,
stream density
\end{IEEEkeywords}

\section{Introduction}

Automatic speech recognition (ASR) on a shared CPU has two service constraints.
Each utterance needs an accurate stable transcript, and the service must keep
up with audio that continues to arrive while neighboring streams are active.
A model that processes a prerecorded file quickly can still fall behind under
periodic arrivals. A system with less average CPU work can consume more memory
per stream or concentrate its work at a chunk boundary. Such effects matter
when a host must serve concurrent captioning, dictation, or voice interfaces
without a GPU.

We study these constraints with the R1 streaming transducer, a CPU inference
engine, and a benchmark stated in deployment units. Final transcript quality is
word error rate (WER). Single-stream cost is CPU-seconds per audio-second
(\cps), measured over an explicitly named boundary. Sustained density is the
count of concurrent paced streams that meet the replay deadline and drain
rules for the tested hardware placement. Figure~\ref{fig:dev-frontier} shows
the cost and quality axes for common LibriSpeech~\cite{panayotov2015librispeech}
development audio. The density table is designed to pair confirmed counts
with held-out quality so that a scheduling gain can be interpreted at its
achieved recognition quality.

We isolate the major changes to R1 through different controls. The dense and
multirate checkpoints test two encoder schedules. A same-weight FP32 comparison
holds the model fixed while changing the runtime. Prepared INT8 is evaluated
against FP32 in the same native engine. Greedy and beam decoding share an
encoder and receive separate rows. The resulting comparisons answer questions
about the implementation, arithmetic, architecture, and decoder at distinct
operating points. An accuracy and cost win against a named control requires
one configuration to improve both metrics on the common workload.

\begin{figure}[!t]
\centering
\input{figure1.tex}
\caption{Pooled word error rate versus single-stream CPU cost on the 730-utterance
LibriSpeech DEV selection panel. \ifpapercoldpoints Each point uses the median
CPU cost of three rotating, single-thread runs on one pinned physical core of
\paperhardware. \else The protocol uses three rotating, single-thread runs on
one pinned physical core of an AMD Ryzen Threadripper Pro 3995WX. \fi Timing
covers the recognizer process tree from launch, including model load, through
the flushed final hypothesis and worker exit. Rust scoring runs afterward.}
\label{fig:dev-frontier}
\end{figure}

\section{Related Work}

Conformer layers combine self-attention with convolution to represent acoustic
dependencies at different scales~\cite{gulati2020conformer}. Transducers model
acoustic and label history jointly and do not require fixed training
alignments~\cite{graves2012transduction}. Their blank transition permits a
decoder to advance in acoustic time, while emitting a label changes the
prediction context. In a streaming implementation, the encoder attention mask,
cache policy, and decoder search govern the actual work required to emit tokens.
The encoder chunk is therefore part of an evaluated configuration.

Reducing the temporal resolution of expensive encoder blocks has precedents
in the temporal U-Net of Squeezeformer and in the multiscale encoder of
Zipformer~\cite{kim2022squeezeformer,yao2024zipformer}. Zipformer also changes
the attention, normalization, and optimization recipes. We use the released
streaming English model in the June 2023 sherpa-ONNX bundle~\cite{zipformermodel}
as the R0 reference. That bundle supplies the encoder, decoder, and joiner as
separate ONNX files. The INT8 model and tokenizer revision are fixed across
the experiments that use R0. Its LibriSpeech training regime makes it a relevant
in-corpus baseline for the research checkpoints, while differences in the training
recipes make system-level WER comparisons descriptive.

Moonshine~v2 uses local sliding-window attention and caches encoder states
as the audio prefix grows~\cite{kudlur2026moonshine}. Its authors evaluate
WER using Transformers on the Open ASR leaderboard. Their separate native
C++ CPU latency experiment on Apple M3 uses a voice-activity detector to mark
the end of speech. In that experiment the authors report 50 and 148~ms
response latency for Tiny and Small respectively, with audio-processing time
equal to 8.03\% and 17.97\% of audio duration. Those numbers describe their
end-of-speech-to-response and processing-time protocol. Our CPU boundary
covers model startup, feature extraction, decode, hypothesis delivery and
shutdown, and our paced experiment measures deadline compliance over the
duration of replay.

Nemotron English and Nemotron 3.5 use cache-aware FastConformer transducers
with selectable chunk settings~\cite{noroozi2024stateful,nemotronen,nemotron35}.
Their C++ inference stack provides quantized GGUF weights and reports the
effective runtime geometry~\cite{nemoasr}. The English Q8 rows use two settings
of one model, so their comparison can reveal sensitivity to right context.
Nemotron 3.5 adds explicit language conditioning. NVIDIA reports parallel
real-time stream counts and median final-token latency on one H100 GPU for
Nemotron 3.5. Those measurements use the vendor GPU and concurrency
protocol. The Nemotron 3.5 CPU adapter passed smoke validation with English
language conditioning in its 80-ms zero-right-context profile.

Low-precision matrix multiplication can reduce weight traffic and use integer
dot products efficiently~\cite{jacob2018quantization}. The benefit depends on
the size of each matrix, the active batch of frame rows, and the overhead of
activation quantization and conversion. The released systems use different
mixed-precision layouts. A model size or a count of nominal integer operations
cannot account for layout, cache reuse, front-end work, or decoder expansions.
Our measured boundary includes all of these costs.

\section{System and Training}
\label{sec:system}

\subsection{Dense and multirate encoding}

Both R1 encoders start from 80-bin log-mel features at 100~Hz. Two causal
convolutional subsampling layers produce 25-Hz encoder states of width 384.
Dense S uses sixteen Conformer blocks with six attention heads. Each block
contains feed-forward, chunk-masked attention, and causal depthwise
convolution operations. The encoder uses rotary positions and a bounded
history. A 160-ms chunk contains four 25-Hz states and attends into the current
chunk and 128 preceding 25-Hz states, where available. The mask excludes keys
after the current chunk. The first position of a chunk can use later positions
within that chunk.

Multirate S\_MR keeps the first four and final two blocks at 25~Hz. A learned
projection maps a pair of consecutive high-rate states into a single 12.5-Hz
state. The middle ten blocks process those lower-rate states with the
corresponding two-state chunk and halved history length. A learned output
projection restores two states per lower-rate state. A high-rate skip carries
the output of the first stack into that reconstruction, and the final stack and
transducer operate at 25~Hz. The design thus keeps the decoder cadence and
the chunk boundary while changing the number of rows processed by the
middle layers. Both rate changes involve trained projections. An odd tail uses
a zero for the missing partner, matching the streaming cache path. This
detail matters because a padded batch row can be nonzero after normalization,
biased projections, and residual blocks.

The dense encoder applies $16 \times 4$ block-frame operations per 160-ms
chunk. Multirate applies $6 \times 4+10 \times 2$, prior to the rate-change
projections. This calculation concerns block inputs. Cost per block also
depends on attention cache lengths, convolution, and the native kernels.
Weight storage for the middle layers persists even when they process fewer
rows, so a lower operation count may give different benefits when one stream
or a batch of streams shares a weight traversal.

\subsection{Training and checkpoint selection}

D2 and D2MR train on the LibriSpeech training partitions totaling about 960
hours. Each started from its corresponding earlier training lineage. D2 uses
pre-normalized residual blocks and an auxiliary CTC objective, with labels
derived from SentencePiece~\cite{kudo2018sentencepiece}. The RNN-T predictor
uses a bounded two-label context. The joiner scores labels and blank from
the current acoustic state and the predictor representation. The auxiliary
head guides training and is independent of the deployed search. The dense
and multirate runs differ in exposure and training history. Their held-out
quality gap measures the available trained systems, so isolating a causal
effect of temporal resolution on WER would require matched recipes and
training budgets.

Selection uses the same speaker-stratified development panel for each
architecture. The candidates include individual checkpoints and adjacent
checkpoint averages. An average combines tensors before export and retains
the topology and parameter shapes of the base checkpoint. This makes it
possible to compare individual and averaged checkpoints with an identical
decoder. The selection criterion is the lowest pooled WER, using the cheaper
configuration to break ties when costs are measured under the same boundary.
We track checkpoint files and export identities independently because a
training checkpoint, an FP32 execution artifact, and a packed INT8 artifact
have different content digests. An evaluated progress checkpoint is labeled
by its step count and remains a progress result while training continues.

\subsection{Runtime and prepared INT8}

The native engine runs feature extraction, encoder attention and convolution,
the prediction network, joiner and detokenization in Rust. Bounded per-stream
key/value and convolution state persists across steps, and scratch storage is
reused across invocations. On the Zen~2 host the native CPU path uses AVX2
and FMA. Packing changes the layout of retained weights. Attention tiling and
cache layouts affect the reuse of acoustic state. The same-weight
ORT-R1 FP32 control compares the exported graph with the native FP32
runtime. Agreement of retained transcripts on the measured workload
checks this control.

Prepared INT8 stores eligible matrices in packed form before the timed
recognizer starts. At run time, the engine scales and centers activations per
frame row for the integer dot product. A row retains its scale when batched
next to a louder or quieter stream. Attention and other operations outside
the integer kernels can remain floating point. Per-row scales constrain the
interactions between streams and make the same input usable in single-stream
and batched tests. Prepared weights also move the artifact conversion out of
the execution lifecycle. The measured boundary still includes reading,
validating and loading the prepared artifact. The same-checkpoint FP32
control establishes the observed numerical-mode trade-off.

\subsection{Decoder operating points}

Greedy RNN-T decoding expands the highest-scoring label until blank advances
time or the symbol cap is reached. Modified beam search maintains several
prefix hypotheses and can recover label choices discarded by greedy search.
A beam result changes the amount and shape of joiner and predictor work,
and can change the time at which a stable word becomes available. We give a
beam transcript its own configuration identity and require an independent
CPU run before pairing that quality with a cost. A decoder comparison on
training evaluation outputs can still quantify accuracy at fixed weights.
It leaves the execution cost of the deployed beam setting open.

\section{Benchmark Protocol}
\label{sec:protocol}

\subsection{Workload and quality}

The frozen DEV selection panel has 400 dev-clean utterances from 40 speakers
and 330 dev-other utterances from 33 speakers. It differs from the older
first-400 dev-clean continuity workload, whose utterances were concentrated
among six speakers. The panel manifest fixes membership and order. The
workload digest binds the exact staged audio and instance identities used by
the recognizer. The benchmark prepares audio and references before the timed
interval. The recognizer receives audio and the driver retains its output.
The Rust scorer normalizes the references and hypotheses after the run. Every
expected utterance must have a retained final result, with all accepted audio
accounted for. The full development sets provide a separate view of checkpoint
and decoder behavior.

For $S$ substitutions, $D$ deletions and $I$ insertions, we report
\begin{equation}
  \operatorname{WER}(U)=100\frac{\sum_{i\in U}(S_i+D_i+I_i)}
                         {\sum_{i\in U}N_i}\,,
  \label{eq:wer}
\end{equation}
where $N_i$ is the normalized reference word count and $U$ is the utterance
population. Pooling sums the numerator and denominator over both splits.
Since split reference counts differ, the mean of the two split percentages
answers a different question. Quality gaps are absolute percentage points
relative to R0. A paired bootstrap resamples speaker clusters, keeping each
recognizer's outputs for an utterance together. For the qualified quality
comparisons, it uses at least 10,000 replicates. Earlier reads of LibriSpeech
test-clean and test-other are part of the project record. Registered sealed
quality results provide the held-out rows, with the prior reads disclosed.

\subsection{Single-stream cost boundary}

For the common DEV comparison we time every system on the same staged
audio, one stream at a time, using one inference thread on one pinned physical
core. We rotate system order across three repetitions to distribute position
effects. Timing starts at recognizer launch and covers all processes in its
tree through final hypothesis flush and worker exit. This includes model
loading, acoustic features, inference, decoder search, output formatting and
drain. Reference preparation, scoring and the CPU use of the observer are
outside the boundary. Each system uses its native streaming interface for audio
and finalization. The driver supplies pre-staged waveform data so that corpus
decoding and staging do not vary by comparator.

The cost statistic is median user plus system CPU seconds per original
audio second across the three complete repetitions. A recognizer running
ahead of real time on unpaced input still accumulates process CPU according
to the work it performs. For each result we retain the per-run receipts,
the sampled peak resident memory of the recognizer tree, and the sizes of
the artifacts that comprise the recognizer. Memory describes loaded and live
state, while artifact size describes stored model data. We record hardware
and build identities, affinity, replay packet timing, decoder, precision, and
effective input step and lookahead. A 160-ms input step and a 160-ms
lookahead can imply different emission delays, even when the same number
appears in the configuration label.

\subsection{Warm cost and paced density}

The separate qualified experiment uses \texttt{ultra bench-suite} on
xbabe2, an AMD Ryzen Threadripper Pro 3995WX host with a benchmark lease.
The in-process warm CPU boundary starts after model load and covers the
work required by a steady-state service. The three native arms in each cost
block are R0, dense R1 and multirate R1. Eight blocks alternate forward and
reverse run order. We calculate the paired log CPU ratio within each block,
report the median ratio with its interval, and apply a one-sided paired
test with a Holm correction across the two comparisons against R0. A
second-day rerun checks whether the median cost agrees within five percent.
Runtime paths whose CPU meters have incompatible boundaries are ineligible
for this warm comparison.

Paced replay feeds 20-ms packets on the capture clock. The shared scheduler
uses a 320-ms step grid with a two-chunk step policy for the internal arms.
Four single-threaded workers share one L3 domain in the per-domain ladder.
The full-host ladder uses 60 workers on 15 cache domains, reserving the
remaining domain for observation and governance. Each ladder starts at
two streams per worker and advances through the preregistered loads. Density
is the measured count for the specified deployment. Scaling the four-worker
result by the number of cache domains has no role in the 60-worker count.

At each tick the scheduler measures the age of the oldest audio packet
awaiting decode against that packet's capture deadline. All workers must
satisfy p99 lateness of at most 800~ms. The final-third mean must stay within
200~ms of the first-third mean, and the run must account for accepted,
processed and unfinished audio at cutoff. These rules distinguish a brief
ability to absorb a queue from keeping up throughout replay. The ladder
uses three 120-s repeats per load and stops at the first failed load.
The last passing load requires all three 300-s confirmations. If confirmation
fails, the confirmed count comes from a lower load that clears its own
confirmation. A pass at the highest tested load is a lower bound. Two
one-hour runs at the claimed density check sustained lateness, backlog and
memory. The accepted and processed audio counters accompany the latency
summary so that incomplete drain is visible.

\subsection{Latency and external adapter validation}

The latency protocol tracks the capture time of each reference event.
Montreal Forced Aligner provides word end times for the
LibriSpeech references~\cite{mcauliffe2017mfa}. A final output is linked
to correctly matched reference words. Word delay uses the first hypothesis
time from which a matched word remains stable through subsequent updates.
We report conditional delay quantiles with coverage, and
\(\mathrm{CQ}(d)\), the fraction of reference words that are correct and
stable within deadline \(d\) of the reference word end. Missing words count
against CQ. With a similar episode denominator, \(\mathrm{FC}(d)\) counts
episodes that finalize within the deadline. We label controlled end-of-input
finalization as such. A deployment with automatic speech endpointing may
show a different finalization delay even with the same encoder.

An adapter must exercise continuous input, retain output timestamps, verify
the end-of-stream flush, and account for processed audio before it enters a
density or latency comparison. A working recorded-file interface can support
the unpaced cost boundary once it has matched staging and successful
drain. For paced replay the driver additionally checks that output is genuinely
produced during replay and that load does not translate into unaccounted
queue growth. External runtime comparators are subject to the same worker
placement and deadline rules. Chunk geometry and right context remain
properties of their respective operating points.

\section{Results}
\label{sec:results}

\subsection{Paced throughput rehearsal}

Table~\ref{tab:rehearsal} shows the passing and first failing rungs in the
same warm bench-suite run on one Zen~2 core complex. Four single-thread
workers replayed the first ten dev-clean utterances in each 20-second rung.
The passing rungs drained all accepted audio. The next rungs exceeded the
800-ms deadline, with worst-worker p99 lateness of
\papernum{rehearsal-r0-pass}{first_fail_p99_ms}\,ms for Zipformer and
\papernum{rehearsal-mr-pass}{first_fail_p99_ms}\,ms for D2MR.

\begin{center}
\begin{minipage}{\columnwidth}\centering
\captionof{table}{Exploratory paced rehearsal on one four-core Zen~2 CCX. Pass and
fail are total streams at the last passing and first failing rungs. The p99
column is worst-worker lateness at the passing rung. WER is pooled over the
730-utterance DEV panel for the same loaded configurations.}
\label{tab:rehearsal}
\begin{tabular}{lrrrr}
\toprule
Configuration & Pass & Fail & p99 (ms) & WER (\%) \\
\midrule
Zipformer INT8 & \metric{rehearsal-r0-pass}{streams_ccx} &
\metric{rehearsal-r0-pass}{first_fail_streams_ccx} &
\metric{rehearsal-r0-pass}{p99_lateness_ms} & \metric{rehearsal-r0-score}{wer_pooled_pct} \\
D2 avg. INT8 & \metric{rehearsal-dense-pass}{streams_ccx} &
\metric{rehearsal-dense-pass}{first_fail_streams_ccx} &
\metric{rehearsal-dense-pass}{p99_lateness_ms} & \metric{rehearsal-dense-score}{wer_pooled_pct} \\
D2MR 50k INT8 & \metric{rehearsal-mr-pass}{streams_ccx} &
\metric{rehearsal-mr-pass}{first_fail_streams_ccx} &
\metric{rehearsal-mr-pass}{p99_lateness_ms} & \metric{rehearsal-mr-score}{wer_pooled_pct} \\
\bottomrule
\end{tabular}
\end{minipage}
\end{center}

At the common 32-stream rung, D2MR used
\pctmetric{rehearsal-mr-common}{cpu_reduction_pct} less CPU per processed
audio second and reduced worst-worker p99 lateness from
\papernum{rehearsal-r0-common}{p99_lateness_ms}\,ms to
\papernum{rehearsal-mr-common}{p99_lateness_ms}\,ms. The dense rehearsal used the
45k/50k/55k average. The DEV-only WER gap between D2MR and Zipformer uses
\papernum{rehearsal-mr-score}{bootstrap_replicates} paired speaker-cluster resamples.
Figure~\ref{fig:rehearsal-load} shows the passing and first failing loads.
Longer runs at the reported load measure sustained capacity.


\begin{figure}[!t]
\centering
\begin{tikzpicture}[x=4.2pt,y=22pt,font=\fontsize{7.7}{9}\selectfont]
  \draw[black!50] (0,0) -- (40,0);
  \foreach \tick in {0,8,16,24,32,40} {
    \draw[black!30] (\tick,0) -- (\tick,3.45);
    \node[below] at (\tick,0) {\tick};
  }
  \foreach \labelname/\id/\y in {Zipformer/rehearsal-r0-pass/1,
                                 Dense D2/rehearsal-dense-pass/2,
                                 D2MR/rehearsal-mr-pass/3} {
    \pgfmathsetmacro{\pass}{\papernum{\id}{streams_per_core}}
    \pgfmathsetmacro{\fail}{\papernum{\id}{first_fail_streams_ccx}/4}
    \fill[black!68] (0,\y-0.13) rectangle (\pass,\y+0.13);
    \draw[dashed] (\pass,\y) -- (\fail,\y);
    \draw (\fail,\y-0.10) -- (\fail,\y+0.10);
    \node[anchor=east] at (-0.25,\y) {\labelname};
  }
  \node[below] at (20,-0.46) {Paced audio streams per physical core};
\end{tikzpicture}
\caption{Measured passing loads in the 20-second replay on one four-core
Zen~2 complex. Dashed segments extend to the first failing load.
Table~\ref{tab:rehearsal} gives the paired DEV WER and passing p99.}
\label{fig:rehearsal-load}
\end{figure}

\begin{table*}[!t]
\caption{Qualified paced density and corresponding test WER. Density columns
count streams for the 4- and 60-worker confirmations. The two WER gaps are
percentage points relative to Zipformer INT8. A dash indicates unavailable
qualified evidence.}
\label{tab:density}
\centering
\small
\begin{tabular}{@{}lrrrrrr@{}}
\toprule
System & 4 workers & 60 workers & Test-clean (\%) & Test-other (\%) &
$\Delta$ clean (pp) & $\Delta$ other (pp)\\
\midrule
\densityrow{R0 Zipformer INT8}{qualified-r0}
\densityrow{Dense R1 prepared INT8}{qualified-dense-r1}
\densityrow{Dense ORT-R1 FP32}{qualified-ort-r1}
\densityrow{Multirate R1 prepared INT8}{qualified-mr-r1}
\densityrow{Moonshine Tiny Streaming}{qualified-moonshine-tiny}
\densityrow{Moonshine Small Streaming}{qualified-moonshine-small}
\densityrow{Nemotron English Q8, 160 ms}{qualified-nemotron-en-160}
\densityrow{Nemotron English Q8, 560 ms}{qualified-nemotron-en-560}
\densityrow{Nemotron 3.5 Q8, English low latency}{qualified-nemotron-35}
\bottomrule
\end{tabular}
\end{table*}

\subsection{Training checkpoint and search controls}

Table~\ref{tab:training} reports full development-set WER for the dense
D2 training checkpoint and its average. The first pair compares greedy
and beam operating points at fixed weights. The two beam cells compare
checkpoint choices at fixed search width.
\paperhasvalue{d2-60k-greedy}{wer_clean_pct}{For the single step
60,000 checkpoint, moving from greedy to width-four beam changes clean WER
from \pctmetric{d2-60k-greedy}{wer_clean_pct} to
\pctmetric{d2-60k-beam4}{wer_clean_pct} and other WER from
\pctmetric{d2-60k-greedy}{wer_other_pct} to
\pctmetric{d2-60k-beam4}{wer_other_pct}. Averaging checkpoints at
50,000, 55,000 and 60,000 further reduces error with beam four to
\pctmetric{d2-avg3-60k-beam4}{wer_clean_pct} clean and
\pctmetric{d2-avg3-60k-beam4}{wer_other_pct} other WER.}{These comparisons
use the trainer evaluation artifacts and comprise accuracy evidence under
that decoder and numerical path.}
We rescore retained trainer hypotheses with the Rust normalizer. The
resulting WERs can differ from the provisional figures produced by the
trainer normalizer. Checkpoint averaging changes weights with the same
topology. The beam changes the search procedure. Its WER reduction
quantifies the search gain at fixed weights and leaves CPU cost to the
matched inference measurement.

\begin{table}[!t]
\caption{Dense D2 checkpoint and search on full development sets. Entries
are split WER (\%). A dash indicates unavailable evidence.}
\label{tab:training}
\centering
\begin{tabular}{@{}lrr@{}}
\toprule
Checkpoint and decoder & Clean & Other\\
\midrule
\trainrow{60k, greedy}{d2-60k-greedy}
\trainrow{60k, beam 4}{d2-60k-beam4}
\trainrow{50/55/60k average, greedy}{d2-avg3-60k-greedy}
\trainrow{50/55/60k average, beam 4}{d2-avg3-60k-beam4}
\bottomrule
\end{tabular}
\end{table}

The full development sets include 2,703 clean and 2,864 other utterances.
The panel pools a deliberately selected subset. A lower WER on the full
sets can coexist with a different ranking on the panel, and the finalist
criterion uses the panel instead.
\paperhasvalue{panel-d2-avg3-60k-greedy}{wer_pooled_pct}{Among the scored
dense checkpoints, the 50/55/60k average is the leading checkpoint. On
the panel, it gives \pctmetric{panel-d2-avg3-60k-greedy}{wer_pooled_pct}
pooled WER with greedy decoding and
\pctmetric{panel-d2-avg3-60k-beam4}{wer_pooled_pct} with width-four beam.
Among the scored D2MR greedy checkpoints, the single 50k snapshot gives
\pctmetric{panel-smr-50k-greedy}{wer_pooled_pct} pooled WER, compared with
\pctmetric{panel-smr-avg3-50k-greedy}{wer_pooled_pct}
and \pctmetric{panel-smr-avg5-50k-greedy}{wer_pooled_pct}
for the three- and five-checkpoint averages. The best observed greedy
points therefore use different checkpoint strategies.}{The tokenizer and
weights of the selected checkpoint carry through export.}
Numerical modes and decoder variants retain distinct configuration IDs.

\paperhasvalue{panel-d2mr-greedy-int8}{int8_delta_pp}{In the engine DEV
check with the D2MR 50k checkpoint, greedy INT8 changes pooled WER by
+\metric{panel-d2mr-greedy-int8}{int8_delta_pp} percentage points against
FP32. The one-sided speaker-cluster upper bound is
+\metric{panel-d2mr-greedy-int8}{int8_upper95_pp} points at 95\% confidence.
The width-four beam check gives a point change of
+\metric{panel-d2mr-beam4-int8}{int8_delta_pp} points and an upper bound of
+\metric{panel-d2mr-beam4-int8}{int8_upper95_pp} points. Both native checks
meet the DEV precision thresholds in the draft protocol. They quantify the
INT8 change at fixed checkpoint and decoder. The associated beam and greedy
INT8 operating points give \pctmetric{panel-d2mr-beam4-int8}{wer_pooled_pct}
and \pctmetric{panel-d2mr-greedy-int8}{wer_pooled_pct} pooled WER respectively.}{}

\subsection{Cold single-stream comparison}

Figure~\ref{fig:dev-frontier} requires eligible CPU measurements
and pooled WER from the same retained hypotheses. The same-weight FP32 pair
speaks to runtime cost, while the FP32 and prepared INT8 rows within the
native engine speak to the arithmetic and artifact change. Dense and
multirate INT8 use the same engine and arithmetic. Their comparison includes
the distinct architecture and training trajectories.
\ifpapercoldpoints
The plotted set consists of the configurations with this complete
timing and transcript evidence.
\else
The common comparison is pending. Training-only WER, unmatched CPU
measurements, and vendor measurements supply context in their stated scopes.
\fi
The streaming cost harness currently supports the greedy R1 path.
Engine and trainer beam evaluations establish accuracy at their checkpoints.
Their streaming cost entries await execution through the common harness.

\subsection{Same-weight development control}

\paperhasvalue{e4-d2-fused-fp32}{cpu_seconds_per_audio_second}{A retained
eight-stream, unpaced CPU experiment on the same 3995WX placement uses
the dense D2 step-5,000 progress checkpoint. The native FP32 runtime
consumed \papernum{e4-d2-fused-fp32}{cpu_seconds_per_audio_second} and ORT FP32
consumed \papernum{e4-d2-ort-fp32}{cpu_seconds_per_audio_second} \cps{} in
their matched run. Their final hypotheses agreed on the continuity
workload. In a later matched build, dense R1 used
\papernum{e4-d2-prepared-int8}{cpu_seconds_per_audio_second} \cps{} in
prepared INT8 and
\papernum{e4-d2-tiled-fp32}{cpu_seconds_per_audio_second} in native FP32.
Zipformer INT8 used
\papernum{e4-zipformer-int8}{cpu_seconds_per_audio_second} \cps{},
with \pctmetric{e4-d2-prepared-int8}{wer_pooled_pct} and
\pctmetric{e4-zipformer-int8}{wer_pooled_pct} WER, respectively. That
operating point saves CPU while increasing errors versus R0.}{The earlier
eight-stream bench isolates the runtime change at fixed FP32 weights and
checks prepared INT8 against a same-runtime FP32 arm. It uses the
step-5,000 dense D2 progress checkpoint and the first-400 DEV continuity
panel. The evidence manifest attaches a separate scope to these records.}
The development control uses a different checkpoint, stream count and
utterance population from Fig.~\ref{fig:dev-frontier}. The historical INT8
runtime also underwent a transcript continuity check across builds. A change
that passes aggregate WER can still change which utterances are correct.
The transcript records therefore bind to the numerical candidate and build.

\subsection{Warm cost, paced service and test accuracy}

Table~\ref{tab:density} provides the main deployment comparison. Each
density entry requires a confirmed count at the tested placement and at the
800-ms p99 deadline. The test-clean and test-other gaps are defined against
the registered Zipformer INT8 result. The bench freeze is still in draft and
the qualified density and test-quality results are pending. Latency coverage
and interval estimates accompany measurements from validated paced adapters
when registered results are available.

Warm cost estimates the CPU work of a loaded system per audio second.
A lower cold cost can result from cheaper startup. Density also depends
on when jobs are runnable and the cost of an iteration with several active
streams. Cache contention, per-stream state, and aligned packet arrivals
can change the capacity of a fully loaded host. These differences motivate
requiring both placement scales at the measured density.

\section{Limitations and Conclusion}

The CPU frontier is conditional on this corpus, host, runtime build and
decoder configuration. LibriSpeech is English read speech. Its dev-clean
and dev-other partitions represent different difficulty within that domain.
They give limited evidence about spontaneous speech, microphone quality,
accent diversity, or background activity. The external models draw on
broader training corpora, with benchmark overlap incompletely identified.
The reported system comparison reflects the weights supplied by each project.
Comparing model design at equal training exposure would require a separate
training experiment.

The cost measurement covers model load through worker exit and is
amortized over the length of the panel. A service that reloads weights for
every connection has a different startup frequency. A persistent service may
hold simultaneous encoder caches for many users and incur scheduling and
memory pressure that the single-stream run cannot exercise. Forced alignment
also contributes timing error to the word-delay metrics, and the use of
controlled end of input leaves production endpointing outside the measurement.
The table's density counts require measured confirmation at each placement.

In the four-core rehearsal, prepared INT8 D2MR reached
\papernum{rehearsal-mr-pass}{streams_per_core} paced streams per core against
\papernum{rehearsal-r0-pass}{streams_per_core} for Zipformer INT8 under the
800-ms p99 deadline. The same configurations differed by
\papernum{rehearsal-mr-score}{wer_pooled_delta_pp} percentage points in pooled
DEV WER. Confirming capacity at that load and at host scale requires the
registered sustained runs. The checkpoint and search comparisons identify
where D2 recovers development accuracy. The evidence manifest links each
operating point to its workload, configuration, and retained outputs.

\bibliographystyle{IEEEtran}
\bibliography{references}
\end{document}
